\createtitlepage{Homework 8}{May 29, 2026}

\begin{problem}[47.5]
  Describe the group of the polynomial $(x^5 - 2) \in (\Q(\zeta))[x]$ over $\Q(\zeta)$, where $\zeta$ is a primitive 5th root of unity.
\end{problem}

\begin{solution}
  Let $\alpha = \sqrt[5]{2}$. We start by finding the degree of the splitting field of $f(x) = x^5 - 2$ over $\Q$. Notice that the roots of $f(x)$ are given by $\alpha, \zeta\alpha, \zeta^2\alpha, \zeta^3\alpha, \zeta^4\alpha$. Since $\zeta$ is a primitive 5th root of unity, the splitting field of $f(x)$ over $\Q$ is $K = \Q(\alpha, \zeta)$. Now, we compute the degree of $[K : \Q]$. Using the tower law, we have:
  \[%
    [K : \Q] = [K : \Q(\zeta)][\Q(\zeta) : \Q]
  .\]%
  We start with $[K : \Q(\zeta)]$. Since $\alpha \not\in \Q(\zeta)$ (as that would mean that $\alpha$ is a root of a polynomial of degree less than 5, which contradicts the irreducibility of $f(x)$), we have $[K : \Q(\zeta)] = 5$. Next, we compute $[\Q(\zeta) : \Q]$. The field $\Q(\zeta)$ is generated by a primitive 5th root of unity, and the minimal polynomial of $\zeta$ over $\Q$ is the 5th cyclotomic polynomial, which has degree $\phi(5) = 4$. Thus, $[\Q(\zeta) : \Q] = 4$. Therefore, we have:
  \[%
    [K : \Q] = 5 \cdot 4 = 20
  .\]%
  Thus, the Galois group of the splitting field of $f(x)$ over $\Q$ is a group of order 20. Therefore, we have $|\Gal(K/\Q(\zeta))| = 20/4 = 5$. Since 5 is prime, there is only one group of order 5, which is $\Z_5$. Thus,
  \[%
    \Gal(K/\Q(\zeta)) \cong \Z_5
  .\qedhere\]%
\end{solution}

\begin{problem}[47.8]
  Find the splitting field $K$ in $\C$ of the polynomial $(x^4 - 4x^2 - 1) \in \Q[x]$. Compute the group of the polynomial over $\Q$ and exhibit the correspondence between the subgroups of $G(K/\Q)$ and the intermediate fields. In other words, do the complete job.
\end{problem}

\begin{solution}
  Let $y = x^2$. Thus, we have $f(x) = x^4 - 4x^2 - 1 = y^2 - 4y - 1$. Solving for $y$, we get:
  \[%
    y = 2 \pm \sqrt{5}
  .\]%
  Thus, the roots of $f(x)$ are given by $\pm\sqrt{2 + \sqrt{5}}$ and $\pm\sqrt{2 - \sqrt{5}}$. Let $\alpha = \sqrt{2 + \sqrt{5}}$ and $\beta = \sqrt{2 - \sqrt{5}}$. We start by finding the degree of the splitting field of $f(x)$ over $\Q$. Computing $\alpha\beta = i$. Thus, we have $\beta = i/\alpha$, meaning $K = \Q(\alpha, i)$. Now, we compute the degree of $[K : \Q]$. Using the tower law, we have:
  \[%
    [K : \Q] = [K : \Q(\alpha)][\Q(\alpha) : \Q]
  .\]%
  We start with $[K : \Q(\alpha)]$. Since $\alpha \in \R$, every element of $\Q(\alpha)$ is also in $\R$. Thus, $i \not\in \Q(\alpha)$. Thus, we have $[K : \Q(\alpha)] = 2$. Next, we compute $[\Q(\alpha) : \Q]$. The field $\Q(\alpha)$ is generated by $\alpha$, and the minimal polynomial of $\alpha$ over $\Q$ is $f(x)$, which has degree 4. Thus, $[\Q(\alpha) : \Q] = 4$. Therefore, we have:
  \[%
    [K : \Q] = 2 \cdot 4 = 8
  .\]%

  There are 5 groups of order 8. The abelian groups of order 8 are $\Z_8$, $\Z_4 \times \Z_2$, and $\Z_2 \times \Z_2 \times \Z_2$ and the non-abelian groups are $D_4$, and $Q_8$. We start by constructing all 8 automorphisms.
  \begin{alignat*}{3}
    \sigma_1 : \alpha &\mapsto \phantom{-}\alpha,\quad&& i \mapsto \phantom{-}i \\
    \sigma_2 : \alpha &\mapsto -\alpha,\quad&& i  \mapsto \phantom{-}i \\
    \sigma_3 : \alpha &\mapsto \phantom{-}i\alpha,\quad&& i \mapsto \phantom{-}i \\
    \sigma_4 : \alpha &\mapsto -i\alpha,\quad&& i \mapsto \phantom{-}i \\
    \sigma_5 : \alpha &\mapsto \phantom{-}\alpha,\quad&& i \mapsto -i \\
    \sigma_6 : \alpha &\mapsto -\alpha,\quad&& i \mapsto -i \\
    \sigma_7 : \alpha &\mapsto \phantom{-}i\alpha,\quad&& i \mapsto -i \\
    \sigma_8 : \alpha &\mapsto -i\alpha,\quad&& i \mapsto -i
  .\end{alignat*}
  We know that $G(K/\Q)$ isn't abelian, since $\sigma_2\sigma_3(\alpha) = -i\alpha$ but $\sigma_3\sigma_2(\alpha) = i\alpha$. Thus, we have $G(K/\Q) \cong D_4$.

  Now, we compute all possible subgroups of $D_4$. Let $\rho = \sigma_3$ and $\mu = \sigma_5$.
  From Lagrange's Theorem, we know that the order of the subgroup must divide 8, meaning it's either 1, 2, 4, or 8. There is only one subgroup of order 1, which is the trivial subgroup, $\{e\}$ and one subgroup of order 8, which is $D_4$ itself. All possible subgroups are given by:
  \begin{gather*}
    H_1 = \braket{\mu}, \quad H_2 = \braket{\rho^2}, \quad H_3 = \braket{\mu\rho}, \quad H_4 = \braket{\mu\rho^2}, H_5 = \braket{\mu\rho^3}, \\
    H_6 = \braket{\rho}, \quad H_7 = \braket{\rho^2, \mu}, \quad H_8 = \braket{\rho^2, \mu\rho^2}
  \end{gather*}

  Now, we compute the fixed fields of each subgroup. Let $\gamma = a + b\alpha + c\alpha^2 + d\alpha^3 + ei + fi\alpha + gi\alpha^2 + hi\alpha^3$.
  \begin{itemize}
    \item $K_{H_1}$: Computing $\mu(\gamma)$, the only unchanged term is $a + b\alpha + c\alpha^2 + d\alpha^3$. Thus, the fixed field $K_{H_1} = \{x \in K | \mu(x) = x\}$ is given by $\Q(\alpha)$.
    \item $K_{H_2}$: Computing $\rho^2(\gamma)$, the only unchanged term is $a + c\alpha^2 + ei + gi\alpha^2$. Thus, the fixed field $K_{H_2} = \{x \in K | \rho^2(x) = x\}$ is given by $\Q(\alpha^2, i)$.
    \item $K_{H_3}$: Computing $\mu\rho(\gamma)$, the only unchanged term is $a + c\alpha^2$. Thus, the fixed field $K_{H_3} = \{x \in K | \mu\rho(x) = x\}$ is given by $\Q(\alpha^2)$.
    \item $K_{H_4}$: Computing $\mu\rho^2(\gamma)$, the only unchanged term is $a + c\alpha^2 - ei - gi\alpha^2$. Thus, the fixed field $K_{H_4} = \{x \in K | \mu\rho^2(x) = x\}$ is given by $\Q(\alpha^2, i)$.
  \end{itemize}
\end{solution}

\begin{problem}[47.9(i)]
  Express the symmetric function $y_1^2 + y_2^2 + y_3^2$ in $y_1$, $y_2$, $y_3$ over $\Q$ as a rational function of the elementary symmetric functions $s_1$, $s_2$, and $s_3$.
\end{problem}

\begin{solution}
  We start by computing $s_1^2$:
  \[%
    s_1^2 = (y_1 + y_2 + y_3)^2 = y_1^2 + y_2^2 + y_3^2 + 2(y_1y_2 + y_1y_3 + y_2y_3) = y_1^2 + y_2^2 + y_3^2 + 2s_2
  .\]%
  Thus, we have $y_1^2 + y_2^2 + y_3^2 = s_1^2 - 2s_2$.
\end{solution}

\begin{problem}[47.10]
  Let $\alpha_1$, $\alpha_2$, and $\alpha_3$ be the zeros in $\C$ of the polynomial
  \[%
    x^3 - 4x^2 + 6x - 2 \in \Q[x]
  .\]%
  Find the polynomial having as zeros precisely the following:
  \begin{enumerate}
    \item $\alpha_1 + \alpha_2 + \alpha_3$.
    \item $\alpha_1^2$, $\alpha_2^2$, and $\alpha_3^2$.
  \end{enumerate}
\end{problem}

\begin{solution}[(i)]
  The polynomial having $\alpha_1 + \alpha_2 + \alpha_3$ as a zero is given by $x - (\alpha_1 + \alpha_2 + \alpha_3)$. We start by computing $\alpha_1 + \alpha_2 + \alpha_3$. By Viète's formulas, we have $\alpha_1 + \alpha_2 + \alpha_3 = 4$. Thus, the polynomial having $\alpha_1 + \alpha_2 + \alpha_3$ as a zero is given by $x - 4$.
\end{solution}

\begin{solution}[(ii)]
  The polynomial having $\alpha_1^2$, $\alpha_2^2$, and $\alpha_3^2$ as zeros is given by $(x - \alpha_1^2)(x - \alpha_2^2)(x - \alpha_3^2)$. We start by computing the elementary symmetric functions of $\alpha_1^2$, $\alpha_2^2$, and $\alpha_3^2$:
  \begin{align*}
    s_1 &= \alpha_1^2 + \alpha_2^2 + \alpha_3^2 = (\alpha_1 + \alpha_2 + \alpha_3)^2 - 2(\alpha_1\alpha_2 + \alpha_1\alpha_3 + \alpha_2\alpha_3) = 4^2 - 2(6) = 16 - 12 = 4, \\
    s_2 &= \alpha_1^2\alpha_2^2 + \alpha_1^2\alpha_3^2 + \alpha_2^2\alpha_3^2 = (\alpha_1\alpha_2)^2 + (\alpha_1\alpha_3)^2 + (\alpha_2\alpha_3)^2 = 6^2 - 4(4)(-2) = 36 + 32 = 68, \\
    s_3 &= \alpha_1^2\alpha_2^2\alpha_3^2 = (\alpha_1\alpha_2\alpha_3)^2 = (-2)^2 = 4.
  .\end{align*}
  Thus, the polynomial having $\alpha_1^2$, $\alpha_2^2$, and $\alpha_3^2$ as zeros is given by $x^3 - s_1x^2 + s_2x - s_3 = x^3 - 4x^2 + 68x - 4$.
\end{solution}

\begin{problem}[47.11]
  Show that every finite group is isomorphic to some Galois group $G(K/F)$ for some finite normal extension $K$ of some field $F$.
\end{problem}

\begin{solution}
  Let $G$ be a finite group. By Cayley's theorem, $G$ is isomorphic to a subgroup of $S_n$ for some $n$, so it suffices to show every subgroup of $S_n$ arises as a Galois group.

  Let $F = \Q$ and let $K$ be the splitting field of $f(x) = \prod_{i=1}^n (x - t_i)$ over the field $\Q(t_1, \ldots, t_n)$ of rational functions in $n$ indeterminates. By a standard result, the Galois group of this generic polynomial is $S_n$. However, it is cleaner to argue directly as follows.

  Let $H \le S_n$ be any subgroup (with $G \cong H$). Let $L = \Q(t_1, \ldots, t_n)$ be a purely transcendental extension of $\Q$, on which $S_n$ acts by permuting the indeterminates $t_i$. Let $F = L^{S_n}$ be the fixed field of $S_n$, which is the field of symmetric rational functions, and let $K = L^H$ be the fixed field of $H$. Since $L/F$ is a finite Galois extension with $G(L/F) \cong S_n$, the Galois correspondence gives
  \[%
    G(L/K) \cong H
  ,\]%
  and $L/K$ is a finite normal extension. Setting $K' = K$ and $F' = L$ (relabeling), we have exhibited $H \cong G(L/K)$ as a Galois group of a finite normal extension. Since $G \cong H$, the group $G$ is isomorphic to the Galois group $G(L/K)$.
\end{solution}

\begin{problem}[48.2]
  Classify the group of the polynomial $(x^{20} - 1) \in \Q[x]$ over $\Q$ according to the Fundamental Theorem of Finitely Generated Abelian Groups. Theorem 9.12. [Hint: Use Theorem 48.4.]
\end{problem}

\begin{solution}
  The splitting field of $x^{20} - 1$ over $\Q$ is the 20th cyclotomic extension $K = \Q(\zeta)$, where $\zeta = e^{2\pi i/20}$ is a primitive 20th root of unity. By Theorem 48.4, the Galois group $G(K/\Q)$ is isomorphic to the group of positive integers less than 20 and relatively prime to 20 under multiplication modulo 20, which is $\Z_{20}^*$, and has order $\phi(20)$.

  We compute $\phi(20) = \phi(4)\phi(5) = 2 \cdot 4 = 8$. The units of $\Z_{20}$ are
  \[%
    \Z_{20}^* = \{1, 3, 7, 9, 11, 13, 17, 19\}
  .\]%
  To classify this group, we use the factorization $20 = 4 \cdot 5$ and the Chinese Remainder Theorem, which gives
  \[%
    \Z_{20}^* \cong \Z_4^* \times \Z_5^*
  .\]%
  Now $\Z_4^* = \{1, 3\} \cong \Z_2$ and $\Z_5^* \cong \Z_4$ (since $\Z_5^*$ is cyclic of order $\phi(5) = 4$). Therefore
  \[%
    G(K/\Q) \cong \Z_{20}^* \cong \Z_2 \times \Z_4
  .\]%
  This is the classification according to the Fundamental Theorem of Finitely Generated Abelian Groups: $G(K/\Q)$ is isomorphic to $\Z_2 \times \Z_4$.
\end{solution}

\begin{problem}[48.5]
  Find the smallest angle of integral degree, that is, $1^\circ, 2^\circ, 3^\circ,$ and so on, constructible with a straightedge and a compass. [Hint: Constructing a $1^\circ$ angle amounts to constructing the regular 360-gon, and so on.]
\end{problem}

\begin{solution}
  Constructing an angle of $k^\circ$ is equivalent to constructing the regular $(360/k)$-gon, which is possible only if $\phi(360/k)$ is a power of 2. We seek the largest $k$ dividing 360 for which this holds, as that gives the smallest constructible integral angle.

  The divisors of 360 corresponding to small angles are $n = 360/k$. By the theory of Section 48, a regular $n$-gon is constructible only if $\phi(n)$ is a power of 2, which requires that every odd prime dividing $n$ is a Fermat prime appearing to the first power. The Fermat primes are $3, 5, 17, 257, 65537$.

  We factor $360 = 2^3 \cdot 3^2 \cdot 5$. For $n \mid 360$, we need $\phi(n)$ to be a power of 2, so $n$ may contain no odd prime to a power greater than 1. Since $3^2 \mid 360$, any $n$ divisible by $9$ fails. Thus $n$ must be of the form $2^a \cdot 3^b \cdot 5^c$ with $b, c \in \{0,1\}$, i.e., $n \in \{1, 2, 3, 4, 5, 6, 10, 12, 15, 20, 24, 60\}$ restricted to divisors of 360 with no squared odd prime factor. The largest such divisor of 360 is $n = 2^3 \cdot 3 \cdot 5 = 120$, giving
  \[%
    \phi(120) = \phi(8)\phi(3)\phi(5) = 4 \cdot 2 \cdot 4 = 32 = 2^5,
  \]%
  which is indeed a power of 2, so the regular 120-gon is constructible. This corresponds to an angle of
  \[%
    k = \frac{360}{120} = 3^\circ.
  \]%
  For $k = 2^\circ$ we would need $n = 180 = 2^2 \cdot 3^2 \cdot 5$, but $3^2 \mid 180$ so $\phi(180) = \phi(4)\phi(9)\phi(5) = 2 \cdot 6 \cdot 4 = 48$, which is not a power of 2, so the regular 180-gon is not constructible. For $k = 1^\circ$ we would need $n = 360 = 2^3 \cdot 3^2 \cdot 5$, and similarly $3^2 \mid 360$ forces $\phi(360) = 96$, not a power of 2.

  Therefore the smallest constructible integral angle is $3^\circ$.
\end{solution}

\begin{problem}[48.8]
  Show that
  \[%
    x^n - 1 = \prod_{d|n} \Phi_d(x)
  ,\]%
  in $\Q[x]$,where the product is over all divisors $d$ of $n$.
\end{problem}

\begin{solution}
  Recall that $\Phi_d(x)$ is the $d$th cyclotomic polynomial, whose roots are exactly the primitive $d$th roots of unity. We show the two sides have the same roots with the same multiplicities, and the same leading coefficient.

  The roots of $x^n - 1$ in $\C$ are exactly the $n$th roots of unity $\{\zeta \in \C \mid \zeta^n = 1\}$, each occurring with multiplicity 1. Every $n$th root of unity $\zeta$ is a primitive $d$th root of unity for exactly one divisor $d \mid n$, namely $d = \text{ord}(\zeta)$, which divides $n$. Conversely, every primitive $d$th root of unity for any $d \mid n$ is an $n$th root of unity. Thus the roots of $\prod_{d \mid n} \Phi_d(x)$ are exactly the primitive $d$th roots of unity ranging over all $d \mid n$, which is precisely the set of all $n$th roots of unity, each appearing exactly once. Hence $x^n - 1$ and $\prod_{d \mid n} \Phi_d(x)$ have the same roots in $\C$, each with multiplicity 1.

  Since both sides are monic polynomials of degree $n$ (as $\sum_{d \mid n} \phi(d) = n$) with the same roots, we conclude
  \[%
    x^n - 1 = \prod_{d \mid n} \Phi_d(x)
  \]%
  in $\C[x]$, and in particular in $\Q[x]$.
\end{solution}

\begin{problem}[48.10]
  Find $\Phi_{12}(x)$ in $\Q[x]$. [Hint: Use Exercises 8 and 9.]
\end{problem}

\begin{solution}
  By Exercise 8 (Problem 48.8), $x^{12} - 1 = \prod_{d \mid 12} \Phi_d(x)$. The divisors of 12 are $1, 2, 3, 4, 6, 12$, so
  \[%
    x^{12} - 1 = \Phi_1(x)\Phi_2(x)\Phi_3(x)\Phi_4(x)\Phi_6(x)\Phi_{12}(x)
  .\]%
  From Exercise 9, the cyclotomic polynomials for small $n$ are:
  \[%
    \Phi_1(x) = x - 1, \quad \Phi_2(x) = x + 1, \quad \Phi_3(x) = x^2 + x + 1, \Phi_4(x) = x^2 + 1, \quad \Phi_6(x) = x^2 - x + 1
  .\]%
  The product of the known factors is
  \[%
    \Phi_1\Phi_2\Phi_3\Phi_4\Phi_6 = (x-1)(x+1)(x^2+x+1)(x^2+1)(x^2-x+1)
  .\]%
  We compute step by step. First, $(x-1)(x+1) = x^2 - 1$. Next,
  \[%
    (x^2-1)(x^2+x+1) = x^4 + x^3 + x^2 - x^2 - x - 1 = x^4 + x^3 - x - 1
  .\]%
  Then,
  \[%
    (x^4 + x^3 - x - 1)(x^2 + 1) = x^6 + x^5 - x^3 - x^2 + x^4 + x^3 - x - 1 = x^6 + x^5 + x^4 - x^2 - x - 1
  .\]%
  Finally,
  \[%
    (x^6 + x^5 + x^4 - x^2 - x - 1)(x^2 - x + 1)
  .\]%
  Expanding:
  \begin{alignat*}{3}
    &x^2(x^6 + x^5 + x^4 - x^2 - x - 1) &&= x^8 + x^7 + x^6 - x^4 - x^3 - x^2, \\
    &{-x}(x^6 + x^5 + x^4 - x^2 - x - 1) &&= -x^7 - x^6 - x^5 + x^3 + x^2 + x, \\
    &1 \cdot (x^6 + x^5 + x^4 - x^2 - x - 1) &&= x^6 + x^5 + x^4 - x^2 - x - 1
  .\end{alignat*}
  Summing by degree:
  \[%
    x^8 + (1-1)x^7 + (1-1+1)x^6 + (0-1+1)x^5 + (-1+1+1)x^4 + (-1+1)x^3 + (-1+1-1)x^2 + (1-1)x - 1
  ,\]%
  which simplifies to $x^8 + x^6 + x^4 - x^2 - 1$. But this has degree 8, whereas $\deg\Phi_{12} = \varphi(12) = 4$ and $\deg(x^{12}-1) = 12$, so we should have $8 + 4 = 12$. Solving for $\Phi_{12}(x)$:
  \[%
    \Phi_{12}(x) = \frac{x^{12}-1}{x^8 + x^6 + x^4 - x^2 - 1}
  .\]%
  Performing polynomial long division of $x^{12} - 1$ by $x^8 + x^6 + x^4 - x^2 - 1$ yields
  \[%
    \Phi_{12}(x) = x^4 - x^2 + 1
  .\qedhere\]%
\end{solution}

\begin{problem}[48.12]
  Let $n, m \in \Z^+$ be relatively prime. Show that the splitting field in $\C$ of $x^{nm} - 1$ over $\Q$ is the same as the splitting field in $\C$ of $(x^n - 1)(x^m - 1)$ over $\Q$.
\end{problem}

\begin{solution}
  Let $\zeta = e^{2\pi i / nm}$, a primitive $nm$th root of unity, so that the splitting field of $x^{nm} - 1$ over $\Q$ is $K = \Q(\zeta)$. Let $L$ be the splitting field of $(x^n - 1)(x^m - 1)$ over $\Q$. The roots of $(x^n-1)(x^m-1)$ are the $n$th and $m$th roots of unity, so $L = \Q(\zeta_n, \zeta_m)$ where $\zeta_n = e^{2\pi i/n}$ and $\zeta_m = e^{2\pi i/m}$ are primitive $n$th and $m$th roots of unity respectively.

  Since $\gcd(n,m) = 1$, by Bezout's theorem there exist $s, t \in \Z$ such that $sn + tm = 1$. Then
  \[%
    \zeta = e^{2\pi i / nm} = e^{2\pi i (sn + tm)/nm} = e^{2\pi i s/m} \cdot e^{2\pi i t/n} = \zeta_m^s \cdot \zeta_n^t \in L.
  \]%
  Hence $K = \Q(\zeta) \subseteq L$.

  The $n$th roots of unity are $nm$th roots of unity (since $(\zeta^m)^n = \zeta^{mn} = 1$), so all roots of $x^n - 1$ lie in $K$. Similarly all roots of $x^m - 1$ lie in $K$. Hence $L \subseteq K$.

  Therefore $K = L$, and the two splitting fields coincide.
\end{solution}
